\subsection{Materia necesaria}

\begin{definicion}{Carta de control para observaciones individuales}
Cuando se tiene una secuencia de mediciones individuales y se asume normalidad, los limites de control se construyen alrededor de la media muestral usando un valor critico \(z\) asociado al nivel de confianza.
\end{definicion}

\begin{formulas}{Media, varianza y limites}
\[
\bar x=\frac{\sum_i x_i}{n},\qquad
\sigma=\sqrt{\frac{\sum_i x_i^2}{n}-\bar x^2}.
\]
\[
LCI=\bar x-z\sigma,\qquad LCS=\bar x+z\sigma.
\]
Para \(96\%\) de confianza se usa \(z\approx 2.055\).
\end{formulas}

\begin{algoritmo}{Control con datos individuales}
\begin{enumerate}
    \item Calcular \(\bar x\) y \(\sigma\).
    \item Construir \(LCI\) y \(LCS\).
    \item Marcar toda muestra fuera del intervalo.
    \item Si se eliminan causas asignables, recalcular limites sin esas muestras.
\end{enumerate}
\end{algoritmo}

\begin{trampa}{Agregados inconsistentes}
En este PDF la tabla transcrita y el agregado impreso \(\sum x_i^2\) no coinciden. En examen, si el enunciado entrega agregados oficiales, se privilegian esos agregados para reproducir la pauta.
\end{trampa}
